\scanpage{17}
% L1–2
Considérons l'hyperalgèbre
\[
\unsure{\mathcal U|U}=A[X]/X^pA[X]
\]
% L3–4
sur l'anneau $A$ \inserted{de car.~$p$}, avec $\Delta$ donné par
\[
\boxed{\Delta X=X\otimes1+1\otimes X}
\]
% L5–8
[\,cela définit à priori $A[X]\mathrel{\unsure{\text{\textemdash}|\longrightarrow}}A[X]\otimes A[X]$, mais
$\Delta X^p=X^p\otimes1+1\otimes X^p$ est \unsure{appliqué|multiple} sur l'idéal
$A[X]\otimes X^pA[X]+X^pA[X]\otimes A[X]$\,], l'augmentation étant définie par
\deleted{\unsure{l'augmentation|l'image}} $\boxed{\varepsilon(X)=0}$, et le $\vee$ par $\boxed{\overset{\vee}{X}=-X}$

% L9–12
Cette hyperalgèbre est isomorphe à sa duale, \deleted{à l'aide}
On la regarde comme \deleted{\unsure{l'anneau des opérateurs|l'algèbre des opérateurs}} dual d'un groupe dans la catégorie des $A$-algèbres commutatives.

% L13–16
Soit $C$ une $A$-algèbre commutative. \unsure{De|Le|La|Se} faire opérer l'hyperalgèbre $\mathcal U$ sur $C$ revient à définir l'image $\deleted{X}\unsure{D'|D^{1}|D}$ de $\deleted{D}X$, i.e. à se donner un opérateur $D$ $A$-linéaire de $C$. De plus
% L17–18
\[
D^p=0\tag*{(i)}
\]
(exprime que $D$ définit une représentation de $A$-algèbre de $A[X]/X^pA[X]$ ds $\unsure{L_A|\mathcal L_A}(C,C)$)
% L19–20
\[
D(xy)=xDy+yDx\tag*{(ii)}
\]
(exprime que cette représentation est compatible avec $\Delta$)

% L21–22
donc \unsure{tt} revient à la \uline{donnée d'une $A$-dérivation} de puissance \unsure{$p^{\text{ième}}$|$p$-ième} nulle.

% L23–27
On veut exprimer maintenant que $C$ est un $\deleted{\unsure{U|H}}\unsure{U'|\mathcal U'|B|B'}$-espace \uline{principal} \uline{homogène} dans la catégorie des $A$-algèbres commutatives. On suppose \deleted{d'abord que} déjà que $C$ \deleted{\unsure{admet une}} \inserted{est \deleted{\unsure{lib|libre}}} un $A$-module libre (c'est sensiblement une
% archive sheet number: 16
